% Ch4 self-study (part 2) -- "I do / You do" ladder for Zumdahl 4.7-4.11.
% Each YOUR TURN carries FOUR questions; bare answers visible in every
% edition (\selfcheck), full solutions key-only. Examples disjoint from
% ch04-notes/worksheets (which own BaSO4, AgCl, and the MnO4-/Fe2+
% titration).
\documentclass{apchem-worksheet}

\apcunitword{Chapter}
\apcunit{4}
\apcunittitle{Types of Chemical Reactions and Solution Stoichiometry}
\apcdoctitle{Self-Study \textbullet\ \S4.7--4.11, I~do / You~do}
\apcsubtitle{Zumdahl \S4.7--4.11 \textbullet\ PDF pp.\ 190--212 \textbullet\ four YOUR TURN questions per ladder}

\begin{document}
\apctitleblock

\begin{apcnote}[How to use these notes --- read this first]
Same ladder rules as part 1: work \textbf{all four} YOUR TURN questions
before checking the gray \emph{check:} line; a tracker tick requires all
four right on the first attempt; full solutions live in the teacher key.

Everything here is Chapter 3 stoichiometry wearing solution clothing:
molarity turns volumes into moles, and from there the balanced equation
takes over exactly as before.
\end{apcnote}

% =====================================================================
\section*{\sffamily Ladder 1 \textbullet\ Precipitation stoichiometry}
\zsec{4.7}

The spine gains one new entrance ramp:
\[ \text{volume} \times M \to \text{moles}
   \xrightarrow{\text{mole ratio}} \text{moles} \to \text{grams} \]
When \emph{two} solution amounts are given, it is a limiting-reactant
problem --- same as Chapter 3.

\begin{workedexample}[I do: lead iodide, the whole problem]
\SI{50.0}{\milli\litre} of \SI{0.200}{\Molar} \ce{Pb(NO3)2} is mixed with
\SI{75.0}{\milli\litre} of \SI{0.100}{\Molar} \ce{NaI}. Find the mass of
precipitate and the concentration of \ce{Pb^2+} left in solution.

\textbf{Net ionic first} (rule 2 exception):
$\ce{Pb^2+ + 2I- -> PbI2(s)}$, $M_{\ce{PbI2}} = \SI{461.0}{\gram\per\mole}$

\textbf{Moles of each reactant:}
\[ \ce{Pb^2+}: 0.0500 \times 0.200 = \SI{0.0100}{\mole} \qquad
   \ce{I-}: 0.0750 \times 0.100 = \SI{0.00750}{\mole} \]

\textbf{Limiting?} The lead present would need $2 \times 0.0100 =
\SI{0.0200}{\mole}$ of iodide; only 0.00750 is there, so \textbf{\ce{I-}
is limiting}.

\textbf{Precipitate} (from the limiting reactant only):
\[ 0.00750 \times \tfrac{1}{2} = \SI{0.00375}{\mole}~\ce{PbI2}
   \;\Rightarrow\; 0.00375 \times 461.0 = \mathbf{1.73~g} \]

\textbf{Leftover \ce{Pb^2+}} --- and here is the solution-specific step:
the total volume is now $50.0 + 75.0 = \SI{125.0}{\milli\litre}$:
\[ n = 0.0100 - 0.00375 = \SI{0.00625}{\mole} \qquad
   [\ce{Pb^2+}] = \frac{0.00625}{0.1250} = \mathbf{0.0500~M} \]
Dividing by the \emph{combined} volume, not the original 50.0 mL, is the
step exams test.
\end{workedexample}

\begin{yourturn}[YOUR TURN 1 --- four questions]
\begin{enumerate}[leftmargin=*,itemsep=0.5em]
  \item Mass of \ce{CaCO3} ($M = \SI{100.09}{\gram\per\mole}$) formed
        when \SI{100.0}{\milli\litre} of \SI{0.150}{\Molar} \ce{Na2CO3}
        reacts with excess \ce{CaCl2}: \workspace[1.8cm]
  \item Mass of \ce{SrSO4} ($M = \SI{183.7}{\gram\per\mole}$) from
        \SI{40.0}{\milli\litre} of \SI{0.250}{\Molar} \ce{Sr(NO3)2} with
        excess \ce{K2SO4}: \workspace[1.8cm]
  \item Volume of \SI{0.100}{\Molar} \ce{AgNO3} required to precipitate
        every chloride ion from \SI{25.0}{\milli\litre} of
        \SI{0.0400}{\Molar} \ce{MgCl2}: \workspace[2cm]
  \item In the worked example, $[\ce{NO3-}]$ after mixing:
        \workspace[1.6cm]
\end{enumerate}
\selfcheck{(a) \SI{1.50}{\gram} \quad (b) \SI{1.84}{\gram} \quad
(c) \SI{20.0}{\milli\litre} \quad (d) \SI{0.160}{\Molar}}
\keyonly{\begin{apcanswer}(a) $0.1000 \times 0.150 = 0.0150$~mol
$\to$ 1:1 $\to$ $0.0150 \times 100.09 = \mathbf{1.50}$~g.
(b) $0.0400 \times 0.250 = 0.0100$~mol $\to$ $0.0100 \times 183.7 =
\mathbf{1.84}$~g. (c) \ce{Cl-}: $0.0250 \times 0.0400 \times 2 =
0.00200$~mol, 1:1 with \ce{Ag+}: $V = 0.00200/0.100 = \mathbf{20.0}$~mL.
(d) Nitrate is a spectator: $n = 2 \times 0.0100 = 0.0200$~mol in
$\SI{125.0}{\milli\litre}$ $\Rightarrow 0.0200/0.1250 =
\mathbf{0.160}$~M.\end{apcanswer}}
\end{yourturn}

% =====================================================================
\section*{\sffamily Ladder 2 \textbullet\ Neutralization stoichiometry}
\zsec{4.8}

Strong acid meets strong base and the net ionic equation is always the
same: $\ce{H+(aq) + OH-(aq) -> H2O(l)}$. The only bookkeeping trap is a
formula that delivers more than one \ce{H+} or \ce{OH-}.

\begin{workedexample}[I do: how much base kills the acid]
What volume of \SI{0.150}{\Molar} \ce{NaOH} neutralizes
\SI{35.0}{\milli\litre} of \SI{0.0900}{\Molar} \ce{H2SO4}?

\textbf{Moles of acid, then of \ce{H+}} --- sulfuric acid is diprotic:
\[ 0.0350 \times 0.0900 = \SI{0.00315}{\mole}~\ce{H2SO4}
   \;\Rightarrow\; 2 \times 0.00315 = \SI{0.00630}{\mole}~\ce{H+} \]

\textbf{\ce{OH-} required} (1:1 with \ce{H+}): \SI{0.00630}{\mole}

\textbf{Volume of base:}
\[ V = \frac{0.00630}{0.150} = \SI{0.0420}{\litre} = \mathbf{42.0~mL} \]

The factor of 2 is the whole exam question. Skip it and every downstream
number is exactly half right --- which is to say, wrong.
\end{workedexample}

\begin{yourturn}[YOUR TURN 2 --- four questions]
\begin{enumerate}[leftmargin=*,itemsep=0.5em]
  \item Volume of \SI{0.100}{\Molar} \ce{HCl} to neutralize
        \SI{20.0}{\milli\litre} of \SI{0.0500}{\Molar} \ce{Ba(OH)2}:
        \workspace[1.8cm]
  \item Molarity of an \ce{HNO3} solution if \SI{25.0}{\milli\litre} of
        it is neutralized by \SI{18.4}{\milli\litre} of
        \SI{0.150}{\Molar} \ce{KOH}: \workspace[1.8cm]
  \item Mass of \ce{Mg(OH)2} ($M = \SI{58.33}{\gram\per\mole}$) needed to
        neutralize \SI{250.}{\milli\litre} of \SI{0.100}{\Molar}
        stomach-acid \ce{HCl}: \workspace[2cm]
  \item Why did question (a) need a factor of 2 while question (b) did
        not? \workspace[1.6cm]
\end{enumerate}
\selfcheck{(a) \SI{20.0}{\milli\litre} \quad (b) \SI{0.110}{\Molar} \quad
(c) \SI{0.729}{\gram} \quad (d) \ce{Ba(OH)2} delivers two \ce{OH-};
\ce{HNO3} and \ce{KOH} are 1:1}
\keyonly{\begin{apcanswer}(a) \ce{OH-}: $0.0200 \times 0.0500 \times 2 =
0.00200$~mol $\Rightarrow V = 0.00200/0.100 = \mathbf{20.0}$~mL.
(b) $n_{\ce{KOH}} = 0.0184 \times 0.150 = 0.00276$~mol $= n_{\ce{HNO3}}$;
$M = 0.00276/0.0250 = \mathbf{0.110}$~M.
(c) $n_{\ce{HCl}} = 0.250 \times 0.100 = 0.0250$~mol; \ce{Mg(OH)2}
neutralizes two each: $0.0125 \times 58.33 = \mathbf{0.729}$~g.
(d) The stoichiometry lives in the formula: \ce{Ba(OH)2} releases two
hydroxides per unit; \ce{HNO3}/\ce{KOH} exchange one proton for one
hydroxide.\end{apcanswer}}
\end{yourturn}

% =====================================================================
\section*{\sffamily Ladder 3 \textbullet\ Titration}
\zsec{4.8}

A titration is a neutralization run \emph{backward}: the known solution
(the titrant) is added until the indicator changes, and the measured
volume reveals the unknown. Vocabulary that gets tested:
\term{equivalence point} (moles of \ce{H+} $=$ moles of \ce{OH-}, set by
stoichiometry) versus \term{endpoint} (where the indicator changes ---
chosen to coincide).

\begin{workedexample}[I do: standardizing a sodium hydroxide solution]
\ce{NaOH} solutions cannot be made to a precise concentration by weighing
(the pellets absorb water), so they are \emph{standardized} against a
weighable acid, potassium hydrogen phthalate (KHP,
$M = \SI{204.22}{\gram\per\mole}$, one acidic H). A \SI{0.6745}{\gram}
KHP sample requires \SI{27.85}{\milli\litre} of the \ce{NaOH} solution.
Find the base's molarity.

\textbf{Moles of acid, from the balance:}
\[ \frac{0.6745}{204.22} = \SI{3.303e-3}{\mole}~\text{KHP} \]

\textbf{At the endpoint, moles of \ce{NaOH} delivered $=$ moles of KHP}
(1:1, one acidic H):
\[ M_{\ce{NaOH}} = \frac{3.303\times10^{-3}}{0.02785}
   = \mathbf{0.1186~M} \]

The mass was weighed to four figures and the volume read to four, so the
answer honestly carries four.
\end{workedexample}

\begin{yourturn}[YOUR TURN 3 --- four questions]
\begin{enumerate}[leftmargin=*,itemsep=0.5em]
  \item \SI{25.00}{\milli\litre} of an \ce{H2SO4} solution requires
        \SI{33.20}{\milli\litre} of \SI{0.1500}{\Molar} \ce{NaOH}. Find
        the acid's molarity: \workspace[2cm]
  \item Vinegar: \SI{25.00}{\milli\litre} requires
        \SI{34.90}{\milli\litre} of \SI{0.1000}{\Molar} \ce{NaOH}. Find
        the mass of acetic acid ($M = \SI{60.05}{\gram\per\mole}$, one
        acidic H) in the sample: \workspace[2cm]
  \item Distinguish the equivalence point from the endpoint in one
        sentence each: \workspace[1.8cm]
  \item Your buret reading at the endpoint was recorded one drop
        ($\approx\SI{0.05}{\milli\litre}$) late. Is the reported acid
        molarity in (a) too high or too low? \workspace[1.6cm]
\end{enumerate}
\selfcheck{(a) \SI{0.0996}{\Molar} \quad (b) \SI{0.2096}{\gram} \quad
(c) stoichiometric equality vs.\ indicator change \quad (d) too high}
\keyonly{\begin{apcanswer}(a) $n_{\ce{NaOH}} = 0.03320 \times 0.1500 =
4.980\times10^{-3}$; diprotic, so $n_{\ce{H2SO4}} = 2.490\times10^{-3}$;
$M = 2.490\times10^{-3}/0.02500 = \mathbf{0.0996}$~M.
(b) $n = 0.03490 \times 0.1000 = 3.490\times10^{-3}$~mol acid $\times\,
60.05 = \mathbf{0.2096}$~g. (c) Equivalence: exactly enough titrant to
consume the analyte, fixed by stoichiometry. Endpoint: where the
indicator visibly changes; a good indicator puts it at the equivalence
point. (d) A late reading overstates the base volume, overstating the
moles of acid --- reported molarity is \textbf{too
high}.\end{apcanswer}}
\end{yourturn}

% =====================================================================
\section*{\sffamily Ladder 4 \textbullet\ Oxidation states; spotting
redox}
\zsec{4.9}

Oxidation states are electron bookkeeping. The rules, in priority order:
element 0; monatomic ion $=$ its charge; F $-1$; O $-2$ (except
peroxides); H $+1$ with nonmetals, $-1$ with metals; everything must sum
to the species' overall charge. Then:

\begin{center}
\fbox{\parbox{0.86\linewidth}{\centering
\textbf{OIL RIG} --- Oxidation Is Loss, Reduction Is Gain (of
electrons). The substance \emph{containing} the oxidized atom is the
\textbf{reducing agent}, and vice versa.}}
\end{center}

\begin{workedexample}[I do: assigning states, then judging a reaction]
\textbf{(a) Chromium in \ce{K2Cr2O7}:} K is $+1$ ($\times 2 = +2$), O is
$-2$ ($\times 7 = -14$); the compound is neutral, so
$2x + 2 - 14 = 0 \Rightarrow x = \mathbf{+6}$.

\textbf{(b) Sulfur in \ce{SO3^2-}:} O contributes $-6$; the ion carries
$2-$, so $x - 6 = -2 \Rightarrow x = \mathbf{+4}$.

\textbf{(c) Is $\ce{Zn(s) + CuSO4(aq) -> ZnSO4(aq) + Cu(s)}$ redox?}
Zn: $0 \to +2$ --- lost electrons, \textbf{oxidized}, so zinc metal is
the \textbf{reducing agent}. Cu: $+2 \to 0$ --- gained electrons,
\textbf{reduced}, so \ce{CuSO4} (the \ce{Cu^2+} in it) is the
\textbf{oxidizing agent}. S and O never change ($+6$, $-2$): sulfate is
along for the ride.
\end{workedexample}

\begin{yourturn}[YOUR TURN 4 --- four questions]
\begin{enumerate}[leftmargin=*,itemsep=0.5em]
  \item Oxidation state of Cl in \ce{ClO3-}: \workspace[1.2cm]
  \item Oxidation state of P in \ce{PO4^3-}: \workspace[1.2cm]
  \item In $\ce{2Al(s) + 3Br2(l) -> 2AlBr3(s)}$, identify what is
        oxidized, what is reduced, and both agents: \workspace[1.8cm]
  \item Is $\ce{CaCO3(s) -> CaO(s) + CO2(g)}$ a redox reaction? Prove
        your answer with oxidation states: \workspace[1.8cm]
\end{enumerate}
\selfcheck{(a) $+5$ \quad (b) $+5$ \quad (c) Al oxidized (reducing
agent); \ce{Br2} reduced (oxidizing agent) \quad (d) no --- no state
changes}
\keyonly{\begin{apcanswer}(a) $x - 6 = -1 \Rightarrow \mathbf{+5}$.
(b) $x - 8 = -3 \Rightarrow \mathbf{+5}$.
(c) Al: $0 \to +3$, oxidized $\Rightarrow$ Al is the reducing agent.
Br: $0 \to -1$, reduced $\Rightarrow$ \ce{Br2} is the oxidizing agent.
(d) \textbf{No.} Ca stays $+2$, C stays $+4$, O stays $-2$ in every
species --- decomposition without electron transfer is not
redox.\end{apcanswer}}
\end{yourturn}

% =====================================================================
\section*{\sffamily Ladder 5 \textbullet\ Balancing redox; a redox
titration}
\zsec{4.10--4.11}

Balancing by \textbf{oxidation states}: find the per-atom electron loss
and gain, choose coefficients so total loss $=$ total gain, then finish
the non-redox atoms by inspection.

\begin{workedexample}[I do: copper dissolving in dilute nitric acid]
\[ \ce{Cu(s) + HNO3(aq) -> Cu(NO3)2(aq) + NO(g) + H2O(l)}
   \qquad\text{(unbalanced)} \]

\textbf{Find the changes:} Cu $0 \to +2$: loses 2. N $+5 \to +2$ (in the
NO only): gains 3.

\textbf{Equalize electrons:} $3 \times \text{Cu}$ loses 6;
$2 \times \text{N}$ gains 6:
\[ \ce{3Cu + HNO3 -> 3Cu(NO3)2 + 2NO + H2O} \]

\textbf{Finish by inspection.} The right side now holds $6 + 2 = 8$ N,
so 8 \ce{HNO3}; 8 H makes 4 \ce{H2O}:
\[ \boxed{\;\ce{3Cu + 8HNO3 -> 3Cu(NO3)2 + 2NO + 4H2O}\;} \]

\textbf{Count everything:} Cu $3=3$; N $8 = 6+2$; H $8 = 8$; O
$24 = 18 + 2 + 4$. Note only 2 of the 8 nitrogens changed state --- the
other 6 became spectator nitrate.
\end{workedexample}

\begin{yourturn}[YOUR TURN 5 --- four questions]
\begin{enumerate}[leftmargin=*,itemsep=0.5em]
  \item Balance the thermite reaction by electron counting:
        $\ce{Al + Fe2O3 -> Al2O3 + Fe}$ \workspace[1.8cm]
  \item Balance by electron counting: $\ce{Mg + N2 -> Mg3N2}$
        \workspace[1.8cm]
  \item A \SI{25.00}{\milli\litre} sample of hydrogen peroxide solution
        requires \SI{31.25}{\milli\litre} of \SI{0.0400}{\Molar}
        \ce{KMnO4}:
        $\ce{2MnO4- + 5H2O2 + 6H+ -> 2Mn^2+ + 5O2 + 8H2O}$.
        Find $[\ce{H2O2}]$: \workspace[2.2cm]
  \item In that titration no indicator is added. What signals the
        endpoint? \workspace[1.6cm]
\end{enumerate}
\selfcheck{(a) \ce{2Al + Fe2O3 -> Al2O3 + 2Fe} \quad
(b) \ce{3Mg + N2 -> Mg3N2} \quad (c) \SI{0.125}{\Molar} \quad
(d) the first drop of permanganate that stays purple}
\keyonly{\begin{apcanswer}(a) Al $0 \to +3$ loses 3; Fe $+3 \to 0$ gains
3: two of each. \ce{2Al + Fe2O3 -> Al2O3 + 2Fe}.
(b) Mg loses 2, N gains 3; equalize at 6: \ce{3Mg + N2 -> Mg3N2}.
(c) $n_{\ce{MnO4-}} = 0.03125 \times 0.0400 = 1.250\times10^{-3}$~mol
$\times \tfrac{5}{2} = 3.125\times10^{-3}$~mol \ce{H2O2};
$M = 3.125\times10^{-3}/0.02500 = \mathbf{0.125}$~M.
(d) While peroxide remains, added \ce{MnO4-} is consumed to nearly
colorless \ce{Mn^2+}; when the analyte is gone, the next drop keeps its
purple color --- permanganate is its own indicator.\end{apcanswer}}
\end{yourturn}

% =====================================================================
\section*{\sffamily Mastery tracker}

Tick a row only if \textbf{all four} YOUR TURN questions were right on
the first attempt.

\renewcommand{\arraystretch}{1.35}
\noindent\begin{tabular}{@{}c p{6.6cm} c p{4.4cm}@{}}
\toprule
\textbf{First try?} & \textbf{Skill} & \textbf{Ladder} &
  \textbf{If not, re-read\ldots}\\
\midrule
$\square$ & precipitation stoichiometry & 1 & combined-volume step \\
$\square$ & neutralization amounts & 2 & the diprotic factor of 2 \\
$\square$ & titration calculations & 3 & KHP standardization \\
$\square$ & oxidation states; redox ID & 4 & the priority rules; OIL RIG \\
$\square$ & balancing redox; redox titration & 5 & loss $=$ gain first \\
\bottomrule
\end{tabular}

\begin{apcnote}[Scoring yourself honestly]
5/5: you have finished Chapter 4's calculations --- the end-of-chapter
problem set is now practice, not new learning. 3--4: redo the missed
ladders tomorrow. 2 or fewer: the recurring root causes here are (1)
molarity $\times$ litres \emph{before} any ratio, (2) the per-formula
count of \ce{H+}, \ce{OH-}, or electrons, and (3) dividing leftover
moles by the \emph{combined} volume. Fix the habit and rerun the
ladders.
\end{apcnote}

\end{document}
